\section{Algorithms and guarantees}
\label{sec:algorithms}

All methods in this section consume the same validated
identity-closure instance.  Strategies and tagged requirements have canonical
original order; weights are exact rationals; mandatory flags, source profiles,
raw module memberships, the source frontier, and the source closure are stored
independently of the incidence matrix.  Preprocessing supplies an exact
objective offset and a reconstruction map.  The number of residual optional
strategies is \(n\), the number of residual tagged requirements is \(r\), and
strategy \(i\) has residual mask \(R_i\) and positive weight \(w_i\).



\subsection{Common output and certificate}
\label{subsec:aor-common-certificate}

Every algorithm returns original strategy identifiers in canonical order, the
exact burden, its declared completion status, its tie rule, an algorithm trace
or search counters, and any required random seed.  A result is reportable only
after an independent certificate recomputes from the original catalog:
\begin{enumerate}
 \item every mandatory strategy is retained;
 \item every tagged requirement is covered;
 \item the complete source operating frontier is preserved exactly;
 \item the complete source identity closure is preserved exactly; and
 \item the exact burden agrees with the returned burden.
\end{enumerate}
Passing this certificate establishes feasibility and burden for the returned
library.  It does not establish heuristic optimality, convert a solver status
into a proof, or show that two implementations are independent.  Complete-tie
output is permitted only when the preprocessing map declares every optimizer
identity reconstructable; otherwise the result is one explicitly declared
representative.

\AORAlgorithmGuaranteeTable

\subsection{Independent exact methods}
\label{subsec:aor-exact-algorithms}

\subsubsection{Complete enumeration}

\begin{proposition}[Correctness of complete enumeration]
\label{prop:aor-complete-enumeration}
Let the input be a valid finite identity-closure tagged instance with exact
nonnegative mandatory weights, strictly positive optional weights, at least one
carrier for every residual requirement, and a valid preprocessing
reconstruction map.  Complete enumeration returns a minimum-burden safe source
sublibrary after checking all \(2^n\) residual selections.  If complete ties
are requested, the map preserves all optimizer identities, and the output cap
is not exceeded, it returns every original optimizer; otherwise it returns the
declared stable representative.
\end{proposition}

\noindent\emph{Proof.}
Every admissible library is a forced set plus one residual subset.
Enumeration evaluates every such subset and takes the exact minimum; the fixed
offset and valid lift preserve burden and feasibility.  The final certificate
maps tagged feasibility back to the original source semantics through
\cref{thm:aor-tagged-cover-equivalence}.  Complete reconstruction follows only
under the stated map and output-cap conditions.  Time is
\(2^n\operatorname{poly}(|I|)\), so enumeration is exact but exponential in
the number of optional strategies.

\subsubsection{Requirement-mask dynamic programming}

Encode a residual requirement subset by an \(r\)-bit mask.  With
\(D_i(q)\) the least exact burden that attains mask \(q\) using the first
\(i\) strategies, the recurrence is
\begin{equation}
 D_i(q)=\min\left\{
 D_{i-1}(q),\
 \min_{p:\,p\mathbin{\mathrm{OR}}R_i=q}
      \bigl(D_{i-1}(p)+w_i\bigr)
 \right\}.
 \label{eq:aor-mask-dp-recurrence}
\end{equation}

\begin{theorem}[Exact requirement-mask dynamic programming]
\label{thm:aor-requirement-mask-dp}
Let the input satisfy every assumption of
\cref{prop:aor-complete-enumeration}.  The recurrence in
\eqref{eq:aor-mask-dp-recurrence} reconstructs at least one minimum-burden safe
source sublibrary.  Representative mode uses \(O(n2^r)\) recurrence operations,
\(O(2^r)\) exact cost states, and \(O(n2^r)\) predecessor records when full
reconstruction is retained.  Value-only two-layer storage uses \(O(2^r)\)
memory.  Materializing \(T\) tied optimizers adds \(O(Tn)\) output bits and may
be exponential in \(n\).  Its bit time is
\(2^r\operatorname{poly}(|I|)\): the method is fixed-parameter tractable in
\(r\), but it is not polynomial in unrestricted input size.
\end{theorem}

\noindent\emph{Proof sketch.}
Induct on \(i\).  At \(i=0\), only the empty subset attains mask zero.  Every
subset of the first \(i\) strategies either skips \(i\), or removes \(i\) to
give a predecessor mask and then takes it.  The two transitions are exhaustive
and construct valid subsets, so the recurrence stores the least exact burden
for each mask.  The all-ones mask is precisely residual feasibility.  Offset,
lifting, and \cref{thm:aor-tagged-cover-equivalence} yield a safe global
optimum.  The state, transition, memory, and output counts give the displayed
bounds.  The complete induction and reconstruction proof is in Online
Resource~1, Section~S5.

\subsection{Weighted greedy construction}
\label{subsec:aor-weighted-greedy}

After selecting every mandatory or forced strategy, let \(U^0\) be the
remaining tagged requirements and let
\[
 d:=\max_i |R_i|,\qquad H(d):=\sum_{j=1}^d\frac1j.
\]

\begin{theorem}[Transferred weighted-greedy guarantee]
\label{thm:aor-weighted-greedy}
Let the source be a finite identity-closure instance satisfying
\cref{thm:aor-tagged-cover-equivalence}; let every residual requirement have a
carrier; let mandatory weights be nonnegative and every optional weight be
positive; and apply only mandatory/forced preprocessing with a valid exact
offset and lift.  If \(U^0\neq\varnothing\), weighted greedy returns a safe
library of burden \(W_G\) satisfying
\[
 W_G\leq H(d)\operatorname{OPT}\leq H(r)\operatorname{OPT}.
\]
Reverse deletion returns a safe, inclusion-wise irreducible library with
\(W_{GR}\leq W_G\), so it satisfies the same bound.  If \(U^0=\varnothing\),
the forced library is optimal and the factor is one.  The guarantee does not
extend to arbitrary nonidentity closure.
\end{theorem}

\noindent\emph{Proof sketch.}
On the residual tagged universe, the rule minimizes weight per newly covered requirement. It is precisely the positive-cost
weighted set-cover greedy algorithm of \citet{Chvatal1979}, which gives
\(G\leq H(d)\operatorname{OPT}_R\).  Adding a nonnegative fixed burden
\(W_0\) preserves the displayed factor.  Reverse deletion accepts only
feasibility-preserving removals of positive-weight strategies and therefore
cannot increase burden.  Later removals cannot restore coverage lost by an
earlier rejected removal, so one reverse pass is irreducible.  The theorem
transfers the classical guarantee through
\cref{thm:aor-tagged-cover-equivalence}; it does not reprove the classical
set-cover theorem.  The cardinality variant replaces \(w_i\) by one in the
score and has no retention-burden approximation guarantee when weights differ.

With straightforward incidence scans, greedy and reverse deletion each use a
polynomial number of exact operations (at most \(n\) selection/deletion rounds
and \(O(nr)\) incidence work per round).  This implementation bound does not
alter the approximation factor.

\subsection{Certified rechecked deletion}
\label{subsec:aor-certified-deletion-suite}

Rechecked deletion starts from the full source and never infers safety from a
stale certificate.  At every scan, each nonmandatory trial deletion is checked
against the original source frontier and closure; after an accepted deletion,
all candidate certificates are discarded and the scan restarts.

Deletion orders include heaviest-first, lightest-first, minimum unique-carrier
exposure, source order, and a seeded random order. Maximum immediate additive
release is the same rule as heaviest-first. The score chooses only among freshly
certified safe deletions; the complete definitions are in Online Resource~1,
Section~S5.

\begin{proposition}[Certified deletion endpoint]
\label{prop:aor-certified-deletion-endpoint}
Let the source be finite and safe relative to itself, let all mandatory
strategies be protected, and suppose every accepted trial deletion is
recomputed exactly and shown to preserve the original frontier, declared
closure, and burden identity.  Then every deterministic or seeded-random
rechecked-deletion variant above terminates at a safe source sublibrary.  A
complete exact final scan certifies that no retained nonmandatory strategy can
be safely deleted.  This is a one-deletion irreducibility certificate, not a
global-optimality or approximation certificate.
\end{proposition}

\noindent\emph{Proof.}
The source is safe.  Exact certification of each accepted trial preserves
safety inductively, and each step removes one strategy, so the trace is finite.
The final scan checks every remaining eligible deletion.  The score and random
order choose only among already certified candidates and therefore do not
change the argument.  If closure is nonidentity, this proposition requires the
separate exact closure checker stipulated in its assumptions.

Multi-start deletion takes the lowest-burden certified endpoint over a fixed,
registered number of random orders. Extra starts cannot worsen the best
observed burden, but supply no approximation factor.

\subsection{Worst-case local--global separation}
\label{sec:aor-safe-deletion-gap}

\begin{theorem}[Unbounded gap for heaviest-safe-first deletion]
\label{thm:aor-heaviest-safe-first-gap}
Let \(k\geq2\) be an integer and let
\(\epsilon\in\mathbb Q\) satisfy \(0<\epsilon<k-1\).  Use one belief,
identity closure, modules \(m_1,\ldots,m_k\), and a mandatory inactive
strategy with zero profile, zero burden, and no modules.  For each \(i\), let
singleton \(s_i\) have zero profile, weight one, and module set \(\{m_i\}\).
Let bundle \(s_B\) have zero profile, weight \(1+\epsilon\), and all \(k\)
modules.  Rechecked heaviest-safe-first deletes the uniquely heaviest safe
bundle, terminates at the \(k\)-singleton library of burden \(k\), and returns
an inclusion-wise irreducible endpoint.  The unique global optimum retains
only the bundle and has burden \(1+\epsilon\).  The exact ratio is
\[
  \frac{k}{1+\epsilon},
\]
so heaviest-safe-first and one-deletion irreducibility alone have no constant
burden guarantee over this family.  The conclusion does not apply to every
deletion order or every certified deletion heuristic.
\end{theorem}

\noindent\emph{Proof.}
All profiles are zero, so safety is exactly coverage of all modules.  At the
source, the bundle makes every singleton deletable and all singletons together
make the bundle deletable.  Since \(\epsilon>0\), the bundle is uniquely
heaviest and is deleted first; afterward each singleton uniquely carries its
module.  Any safe library without the bundle must therefore retain all
singletons and costs \(k\), whereas the bundle alone is safe and costs
\(1+\epsilon<k\); positive weights make it the unique optimum.  Fixing, for
example, \(\epsilon=1/2\) makes the ratio unbounded as \(k\) grows.  A
singleton-first order instead reaches the bundle optimum, proving the stated
order boundary.  The complete family proof and finite fixture are in Online Resource~1, Section~S6.

\subsection{Guarantee and evidence boundary}

Enumeration and mask DP are exact finite methods under their stated
identity-closure assumptions; the former is exponential in \(n\), and the
latter is exponential in \(r\).  Weighted greedy is approximate but has the
transferred \(H(d)\) bound; reverse deletion preserves that bound.  Rechecked
deletion and multi-start return exactly certified feasible, irreducible
endpoints but remain heuristics.  The MIP in
\cref{subsec:aor-baseline-mip} supplies solver evidence plus an exact
returned-library postcheck.  Complete proofs and machine-readable certificate
schemas are in Online Resource~1; benchmark runtimes and method comparisons
are experimental evidence reported separately in Section~\ref{sec:computational-study}.

The negative guarantee concerns this named rule and one-deletion neighborhoods; it makes no claim about stronger exchange search.
